\section{Verification design and limitations}

The proofs establish the all-$n$ claims without enumeration. The
accompanying standalone author verifier enumerates every
member of $\On$ for $1\le n\le7$, a total of $2353$ maps. It constructs
literal successors and whole-function orbits, then compares them with
the image condition, terminal/entrance formulas and sharp witnesses.
For every labelled target, it compares the complete observed predecessor
list with the gap decoder and the Laurent count. Its full output records
all sources and targets, including zero fibres, with explicit canonical
ordering. Such finite checks pressure the proofs but cannot prove their
all-size statements.

The new verifier has completed one initial production and a separately
recorded strict author replay pair. Each invocation checked $14523$
predicates over the same $2353$ maps. The complete initial stdout was
adopted as the canonical; three native comparisons checked each replay
against that canonical and the two replays against one another. These
are supporting finite checks, not additional proofs or independent
manuscript reviews; the execution milestone supplies neither a manuscript
build nor review acceptance. An earlier exploratory author pilot remains
separate: its strict runtime-prelock audit failed and its archived output
is not reused as the canonical or as strict manuscript verification.
The new initial production and pair have separately recorded source,
parameter and bounded runtime-closure evidence.

The scope is the exact clock and one-step inverse of
\eqref{eq:update}. We do not determine a global maximum fibre or its
maximizers, basin cardinalities, or all-time inverse sets. The cited
support/rank structure, static projection criteria and generic deletion
are not new results. Our comparison with those sources is bounded to
their stated objects; it is not a proof excluding every possible owner,
factor, lift or transfer. The short note supplies explicit formulas for
this recomputed operation, not a broader classification of finite
semigroup feedback systems.
